\documentclass[11pt]{article}

% ── Packages ────────────────────────────────────────────────
\usepackage[margin=1in]{geometry}
\usepackage{amsmath, amssymb, amsthm}
\usepackage{booktabs}
\usepackage{siunitx}
\usepackage{graphicx}
\usepackage{hyperref}
\usepackage{natbib}
\bibliographystyle{plainnat}
\usepackage{setspace}
\usepackage{microtype}
\usepackage{enumitem}
\usepackage{caption}
\usepackage{subcaption}
\usepackage{xcolor}
\usepackage{float}
% code packages
\usepackage{listings}
\usepackage{inconsolata}
\lstset{basicstyle=\ttfamily\footnotesize,breaklines=true}
\usepackage{csvsimple}

% ── Hyperref setup ──────────────────────────────────────────
\hypersetup{
    colorlinks = true,
    linkcolor  = black,
    citecolor  = black,
    urlcolor   = blue
}

% ── Theorem environments ────────────────────────────────────
\newtheorem{theorem}{Theorem}
\newtheorem{proposition}{Proposition}
\newtheorem{lemma}{Lemma}
\newtheorem{corollary}{Corollary}
\theoremstyle{definition}
\newtheorem{definition}{Definition}
\newtheorem{assumption}{Assumption}
\theoremstyle{remark}
\newtheorem{remark}{Remark}

% ── Spacing ─────────────────────────────────────────────────
\onehalfspacing

% ── Title block ─────────────────────────────────────────────
\title{The Role of Product Variety in Targeted Pricing}
\author{Tate Mason}
\date{\today}

% ============================================================
\begin{document}
\maketitle
\thispagestyle{empty}
\newpage

%============================================================%
\section{Introduction}
%============================================================%

%============================================================%
\section{Related Literature}
%============================================================%

% LOV Demand Literature

The variety-seeking and state dependence literature has been primarily focused on understanding the dynamics of consumer preferences over time. Earlier papers in this literature, like those of \citet{mcalister_dynamic_2026} and \citet{lattin_model_1987} are most closely linked with this paper's interpretation of variety. They posit that a consumer's preferences for a specific product attribute is satiated with repeated purchases. McAlister finds her model of dynamic attribute satiation outperforms models with independent choice occasions in accuracy and significance of estimates. Lattin finds that there is heterogeneity in the variety-seeking types, though he uses a tiny cross-section of his own students for the study. The satiation-based approach was then built upon by \citet{givon_variety_1984} who used an inertia model to show how consumers can differ by preference for brand versus preference for variety. By using the brand rather than a specific attribute of a product, Givon is not able to pick up on specific drivers of product switching. \citet{erdem_decision-making_1996} also looks at the brand choice problem, modeling decision making when consumers do not have perfect knowledge of preferences over products. The authors use a Bayesian updating framework by which consumers can learn about products via experience or advertising. Their key finding is that consumers are risk averse in consumption, leading them to stay with brands they are familiar with. This model is informative as the reverse-side of my supply-side framework. Marketing researchers, \citet{chintagunta_variety_1999} and \citet{dube_multiple_2004}, greatly inform the demand side of this paper. \citet{chintagunta_variety_1999} shows the validity of what he terms the "Lightning Bolt" model of brand choice. He allows for time-varying variety seeking, implying the variety-seeking type can change with time. He finds that state dependence over time and marketing variables are informative in demand estimation. \citet{dube_multiple_2004} progresses Chintagunta's work, showing how consumers seek variety in quantities purchased through the lens of carbonated sodas. Dub\'e successfully estimates consumer preferences over bundles of products within a category, finding that consumer demographics were key in determining the joint distribution. Both of these papers are primarily concerned with brand choice rather than classifying types of consumers. Finally, \citet{ackerberg_advertising_2003} adds to Erdem and Keane's work, allowing for consumers to learn about product attributes via advertising. The key finding here is that switching is primarily induced by the signals put out by advertisements.

% Supply Literature

The literature on personalized pricing is mostly theoretical for physical goods. \citet{grossman_informative_1984} develops a model of advertising expenditure in which the firm is trying to inform consumers of their products' attributes. Their key finding is that advertising provides a social function by improving consumer and product matching. They also find there are excessive levels of advertising in nearly all cases. \citet{karle_imperfect_nodate} models optimal advertising strategies for the firm. Using a theoretical model with two firms, he finds that firms' profits change non-monotonically with targeting ability, decreasing as targeting improves before increasing again. \citet{moorthy_targeting_2023} derives the optimal conditions for targeted advertising. Using a Hotelling framework, the authors find that the key driver of optimal expenditure levels is product differentiation, claiming it as a managerial framework in the face of targeting regulation. \citet{dube_personalized_2023} uses a randomized control trial (RCT) to understand how price personalization impacts firm profits and consumer welfare. He finds that total consumer surplus declines with personalized pricing, 60\% of consumers benefit and, under some conditions, consumer welfare may increase.

% Where I Fit

My contribution is twofold. Firstly, I am primarily concerned with attribute switching rather than brand switching. Much of the previous work focuses on brand level switching. When looking at variety in attributes, I am able to see how consumers switch both within brand as well as between brands. On the supply side, I develop a simple model which can derive optimal pricing depending on consumer switching-type. This model is then used to simulate the effect of policy around targeted advertising by obscuring the firm's information around consumer types. 

%============================================================%
\section{Data}
%============================================================%

    The data comes from Nielsen's HomeScan and MarketScan. These two main datasets give me access to household purchase decisions as well as product characteristics. From HomeScan, I can access household characteristics, purchase histories by shopping trip, and information on purchase characteristics such as discount use or the total price paid. MarketScan provides product characteristics like price or if the product was in a featured position on the shelf. I limit the sample to Atlanta, Chicago, Denver, Houston, Des Moines, Phoenix, Philadelphia, and San Diego for years 2022-2024. These markets were chosen since they are major cities in different parts of the country. I only include single-person households to avoid any contamination of shopping for multiple people like children or spouses. This ensures I am picking up true preferences for variety rather than buying for multiple people. These trips are encoded with unique identifiers, thus allowing me to look at per-trip, and within-trip, decisions. Agents are under observation for eight weeks and asked to record each shopping trip. Stores are restricted to be valid Nielsen stores, meaning they appear in the reatil code data. This is necessary to observe product prices.

    I use yogurt data as the product space for my analysis. Yogurt was chosen due to its relative lack of storability, wide menu of flavors, as well as diverse nutrition characteristics. I classify flavors as berry, other flavors, and plain. Plain is a unique identifier while "berry" and "other flavors" are aggregates of many flavors. Berry includes any flavor with "berry" in its flavor encoding, as well as cherry flavors. Other flavors are composed of all flavors which are neither plain nor other. By construction, it is the largest category. The outside option is defined as any purchase which is not yogurt. 

%============================================================%
\section{Demand Model}
%============================================================%
    In each period $t$, the consumer $i$ chooses a product $j$ from the menu of available products $J_t$. The available products are $J_t = \{0,1,2,3\}$ where $j=0$ is the outside option, $j=1$ is purchasing berry yogurt, $j=2$ is purchasing plain yogurt, and $j=3$ is purchasing other flavors. Consumers are assumed to be state-dependent, carrying previous choices forward when love of variety is present. Otherwise, they only internalize characteristic utiltiy and price utility. Each product $j\in J_t$ is a bundle of $k\in K$ characteristics such that $X_{jt}^k \in \mathcal{X}_{jt}$. The consumer's utility from choosing product $j$ at time $t$ is given by:


    \begin{alignat*}{1}
      U_{ijt} &= \beta_0 + \beta_{it}X_{jt}
                  + \gamma_i\Xi_{ijt} + \alpha_i p_{jt} +  \varepsilon_{ijt}, \ \text{for j$\in\{1,2,3\}$} \\
    \text{such that} \\
      \Xi_{ijt} &= \lambda_i\sum_{s=1}^{8}\delta^{t-s}\Xi_{ijs} + (1-\lambda_i)\sum_{j\ne k}\rho x_{k,t-1} \\
      \varepsilon_{ijt} &\sim \text{T1EV}(0,1) \\
      \text{$U_{i0t}$}&=\varepsilon_0 
    \end{alignat*}

    In the utility function, there are the following components: $\beta_0$ is a constant, $\beta_{it}$ is a consumer's utility for characteristic $X_{jt}$ associated with product $j$ in time $t$. Following \citet{lattin_model_1987}, love of variety preference is captured by $\gamma_i$. A consumer's history is given by $\Xi_{ijt}$ where $\lambda_i$ is a weight on state dependence $\Xi_{ijs}$ which is discounted with time $\delta^{t-s}$. $\rho$ is an indicator for if a consumer bought a product $x_{k,t-1}$ where $k\ne j$. $\alpha_i$ is the parameter on prices $p_{jt}$ and $\varepsilon_{ijt}$ is the Gumbel shock.

    Choice probabilities take the following form:

    \begin{alignat*}{3}
      P_{ijt} = \frac{\exp(\beta_0 + \beta_{it}X_{jt} + \gamma_i\Xi_{ijt} + \alpha_ip_{jt} +  \varepsilon_{ijt})}{1 + \sum_k\exp(\beta_0+\beta_{it}X_{kt} + \gamma_i\Xi_{ikt} + \alpha_ip_{kt} + \varepsilon_{ikt})}
    \end{alignat*}

    Where we normalize the outside option, not purchasing yogurt, to be zero by construction. 

%============================================================%
\section{Estimation}
%============================================================%

Following \citet{petrin_control_2010}, I use a control function to control for price endogeneity. I do a first stage instrumental variable estimation using Hausman instruments and then take residual from that estimation and insert it into the utility function. The first stage takes the form:

    \begin{alignat*}{2}
      p_{jt} = \bar{p}_{j,-m(i),t} + \tau_t + \eta_{jt}
    \end{alignat*}

    Where $\bar{p}_{j,-m(i),t}$ is the mean price from the three markets the household is not shopping in and $\tau_t$ is a time indicator. The mean price of the other three markets was chosen as an instrument to separate any region specific effects like local advertising. Thus, I assume there is no national advertising or other effects which impact demand across all markets. The residual recovered from the first stage corrects for the bias present in pricing, allowing us a consistent estimate of pricing.

I make an index of prices, instrumental variable residuals, and previous choices for each household's trips for each flavor. These are then used to estimate the objective function:

\begin{alignat*}{4}
  U^k_{ijt} = \beta_0 +  \sum_k\beta_i^kX^k_{jt} + \gamma_i\Xi_{ijt} + \alpha p_{jt} + \sigma\hat{\eta}_{jt}
\end{alignat*}

Where $\beta_0$ is a constant, $\beta^k \ \text{for k} \ \in \{\text{berry, other, plain}\}$ are choice specific coefficients for choosing berry, other, or plain. $\gamma$ is the love of variety term, interacted with $\Xi_{ijt}$, which is just capturing if consumer $i$ switched between periods. The price and control function term are unchanged from above. The probability of individual $i$ choosing the observed product $j$ in time $t$ is given by:
\begin{alignat*}{5}
  \Pi_i(P_{ijt}(\beta_it, \Xi_{ijt}))^{y_{ijt}}
\end{alignat*}

Where $y_{ijt} = 1$ when person $i$ chosen $j$ that period, and zero otherwise. The log-likelihood is then given by:
\begin{alignat*}{6}
  \text{LL}(\theta) = \sum_t\sum_i\sum_jy_{ijt}\ln P_{ijt}(\beta_i, \Xi_{ijt})
\end{alignat*}

I also repeat estimation under specifications which exclude $\gamma_i$, do not estimate heterogeneous $\beta_i, \ \gamma_i$, and do both. This allows me to see if results are biased when love of variety is excluded from estimation. 

%============================================================%
\section{Demand Results}
%============================================================%

Table \ref{satiation_specifications} shows the main results from the estimation. All coefficients are flavor-specific and are relative to the "other flavors" category which is relative to the outside option. When looking at the standard logit, without heterogeneous consumers, I find that consumers experience disutility from both berry flavored and plain yogurt compared to other flavors. They are experience disutility from variety, with berry inducing more disutilty than plain, though plain is not statistically significant. Price effects are negative and significant, with the control function residual being positive and significant. Now, turning to the mixed logit, I find similar estimates. I recover statistically significant coefficients for all but the coefficient on switching for berry yogurt. When looking at the coefficient on plain, I find a large and significant disutility for switching from plain yogurt. I also find higher disutility for consuming plain yogurt. 

Looking at table \ref{tab:naive_specifications}, the estimates for what I term the "naive" specifications. Here, there is no love of variety $\gamma$. When 


%============================================================%
\section{Supply Side}
%============================================================%

There are $N$ consumers, each endowed with a type $\gamma_r$ for $r\in\{H,L\}$, corresponding to high or low utility from variety in products consumed. As such, there are $N_r\in N$ consumers of each type, $r\in\{H,L\}$. Consumers maximize utility according to the following:

\begin{gather*}
  U_{irjt} = \beta_0 + \gamma_r\ind\{j_t\ne j_{t-1}\} + \alpha p_{rjt} +
  \varepsilon_{irjt}
\end{gather*}

Where utility is over individuals $i\in N$, type $r\in\{H,L\}$, products $j\in J$, and time $t\in T$. $\beta_0$ is an intercept, representing the baseline utility from consuming. $\gamma_i$ is the consumer's utility from consuming different products, and is active if they switched products from one period to the next. $\alpha$ is disutility of price, which can be lowered via a discount from advertising. Finally, $\varepsilon_{ijt}$ is the Gumbel error, giving us a closed form Logit.

This equation can be written in terms of deterministic vs. random utility:

\begin{gather*}
  U_{irjt} = V_{irjt} + \varepsilon_{irjt} \\
 \text{s.t.} \\
  V_{irjt} = \beta_0 + \gamma_r\ind\{j_t\ne j_{t-1}\} + \alpha(p_{jt} - d_{rjt})
\end{gather*}

Choice probabilities of consuming the inside option(s), as well as the outside option utility normalized to be zero, is written as such:

\begin{gather*}
  P_{irjt} = \frac{e^{V_{irjt}}}{1 + \sum_ke^{V_{irkt}}}\\
  P_{ir0t} = \frac{1}{1+\sum_ke^{V_{irkt}}}
\end{gather*}

Next, because consumers of the same type $r$ are homogenous, we can say:

\begin{gather*}
  s_{rjt} = \sum_iP_{irjt} \ \text{for r $\in$ \{H,L\}}
\end{gather*}

Taking the derivative with respect to prices, we can get the own and cross price
elasticities:

\begin{gather*}
  \epsilon_{jj}^p = \alpha p_{jt}(1-s_{rjt}) \\
  \epsilon_{jk}^p = -\alpha p_{kt}s_{rkt} \ \ \text{for
  k$\ne$j}
\end{gather*}

Thus, since $\alpha<0$, we can see that own price elasticity is negative and
cross-price is positive. 

Now, a supermarket enters the market. The supermarket is able to control prices
$\bar{p}_j$ for a product $j$. The supermarket can also offer a discount
$d_{rjt}$ which is tailored to our switching types $r\in\{H,L\}$. The firm's
actual price is written as:

\begin{gather*}
  p_{rjt} \equiv \bar{p}_j - d_{rjt} \\
  \text{s.t.} \\
  d_{rjt}>0
\end{gather*}

So, as the market raises the discount, utility increases as $\alpha<0$ so $-\alpha d_{rjt}>0$. It is through this discounting that firm can induce switching. The firm's profit margin for a product j is given by

\begin{gather*}
  m_{rjt} = p_{rjt} - c_j \\
  = \bar{p}_j - d_{rjt} - c_j
\end{gather*}

This says that as we raise the discount, we raise market share on the product being discounted but lose profit margin. 

With these definitions in mind, we will be looking at a two period model under two specifications. First, the case where the firm has perfect information on the consumers' types and can thus perfectly price. In the second case, we will look at the case of imperfect information, where the firm is using Bayesian updating to approximate the type distribution.

\subsection{Perfect Information}

Assume the firm has perfect information over the distribution of types. That is, the supermarket perfectly observes each consumer type $r$ and thus knows the mass $N_r$ of each type. Define period one profit, $\pi_1$ as such:

\begin{gather*}
  \pi_1 = \sum_rN_r\sum_js_{rj1}(m_{rj1}) \\
  \text{s.t.} \\
  s_{rj1} = P_{irj1} \\
  m_{rj1} = \bar{p}_j - d_{rj1} - c_j
\end{gather*}

Here, the firm is understanding how each type's share for each product impacts their profits. We can similarly define period two profits, $\pi_2$ like so:

\begin{gather*}
  \pi_2 = \sum_rN_r\sum_js_{rj1}[\sum_ks_{rk2}(j)(m_{rk2}(j))]
\end{gather*}

The only augmentation is the firm now understands all of the possible paths for a consumer of type $r$ who chose product $j$ in period 1, and product $k$ in period 2. For the sake of conciseness in notation, $w_r(j) = \sum_ks_{rk2}(j)(m_{rk2}(j))$. Define total profit, $\Pi$, as such:

\begin{gather*}
  \Pi = \pi_1 + \pi_2 \\
  = \sum_rN_r\sum_js_{rj1}(m_{rj1} + \delta w_r(j))
\end{gather*}

Where $\delta$ is a discount factor on future profits. Taking the derivative with respect to period one prices for a product $j$ yields:
\begin{gather*}
  \frac{\partial \Pi}{\partial p_{rj1}}: 0 = 1 + \alpha\sum_{k\ne j}s_{rk1}[(m_{rj1} - m_{rk1}) + \delta(w_r(j) - w_r(k))]
\end{gather*}

The firm sets prices according to the current difference in profit margins and the difference in future profits which occur upon switching a consumer between product $j$ and product $k$. When we do the same for period two, we recover the following:

\begin{gather*}
  \frac{\partial \Pi}{\partial p_{rk2}}: 0= 1+\alpha\sum_{l\ne k}s_{rl2}(j)[(m_{rk2}(j) - m_{rl2}(j))]
\end{gather*}

Now, expanding the expression and solving for period one prices:

\begin{gather*}
  0 = 1 + \alpha[\sum_{k\ne j}s_{rk1}((p_{rj1} - c_j) - (p_{rk1} - c_k) + \delta(w_r(j) - w_r(k)))] \\
  0 = 1 + \alpha(p_{rj1} - c_j)\sum_{k\ne j}s_{rk1} - \alpha\sum_{k\ne j}s_{rk1}(p_{rk1} - c_k) + \alpha\delta\sum_{k\ne j}s_{rk1}(w_r(j) - w_r(k)) \\
  \alpha(p_{rj1} - c_j)\sum_{k\ne j}s_{rk1} = -1 + \alpha\sum_{k\ne j}s_{rk1}(p_{rk1} - c_k) - \alpha\delta\sum_{k\ne j}s_{rk1}(w_r(j) - w_r(k)) \\
  \alpha p_{rj1}\sum_{k\ne j}s_{rk1} = \alpha c_j\sum_{k\ne j}s_{rk1} - 1 + \alpha\sum_{k\ne j}s_{rk1}(p_{rk1} - c_k) - \alpha\delta\sum_{k\ne j}s_{rk1}(w_r(j) - w_r(k)) \\
  p^*_{rj1} = c_j - \frac{1}{\alpha\sum_{k\ne j}s_{rk1}} + \frac{\sum_{k\ne j}s_{rk1}(p_{rk1} - c_k)}{\sum_{k\ne j}s_{rk1}} - \delta\frac{\sum_{k\ne j}s_{rk1}(w_r(j) - w_r(k))}{\sum_{k\ne j}s_{rk1}} \\
  p^*_{rj1} = c_j - \frac{1}{\alpha(1-s_{rj1})} + \frac{\sum_{k\ne j}s_{rk1}m_{rk1}}{1-s_{rj1}} - \delta\frac{\sum_{k\ne j}s_{rk1}(w_r(j) - w_r(k))}{1-s_{rj1}} \\
\end{gather*}

From this, the firm with perfect information sets each type $r$'s price to equal the cost plus a demand markup and the share weighted average margin of alternatives $k\ne j$ minus the intertemporal tradeoff. Here, the firm will be willing to lower the price of good $j$ today if it implies higher profits tomorrow ($w_r(j) > w_r(k)$), increasing probability of purchase. In the opposite case, the firm would raise prices on product $j$ to encourage shifting from product $j$ to product $k$ tomorrow. That is, they use the past choice to determine consumer state and thus future profits for each product.
The firm again thinks of the tradeoff, without having to think about future profit as period two is terminal. Under perfect information, the supermarket chooses its prices by evaluating the current margin and future profit associated with reallocating consumers across products. Now, we move to the case of imperfect information.

\subsection{Imperfect Information}

In this case, the supermarket does not know consumer types with certainty, nor do they know the mass of either type. So, they must infer the distribution from from consumer purchases. Define the following objects:
\begin{gather*}
  \lambda_1 = Pr(r = H) \\
  \lambda_2 = Pr(r=H|j) \Rightarrow \lambda_2(j) = \frac{\lambda_1s_{Hj1}}{\lambda_1s_{Hj1} + (1-\lambda_1)s_{Lj1}} \\
  \tilde{s}_{j1} = \lambda_1s_{Hj1} + (1-\lambda_1)s_{Lj1} \\
  \tilde{s}_{k2}(j) = \lambda_2(j)s_{Hk2}(j) + (1-\lambda_2(j))s_{Lk2}(j)
\end{gather*}

Now, we can define the profit functions much like we did before:
\begin{gather*}
  \pi_1 = N\sum_j\tilde{s}_{j1}m_{j1} \\
  \pi_2 = N\sum_j\tilde{s}_{j1}\sum_k\tilde{s}_{k2}(j)m_{k2}(j)
\end{gather*}

The key difference here is that the firm must now update their beliefs from period to period rather than being able to discount according to the tradeoffs between two products. This implies that there is room for error in discounting. When combining the two into overall profit, we get:

\begin{gather*}
  \Pi = N\sum_j\tilde{s}_{j1}[m_{j1} + \delta\sum_k\tilde{s}_{k2}(j)m_{k2}(j)]
\end{gather*}
Now, we can take the derivative with respect to period one prices for product $j$:

\begin{align*}
  \frac{\partial \Pi}{\partial p_{j1}}: 0 &= N\tilde{s}_{j1}+\alpha N[\lambda_1s_{Hj1}(1-s_{Hj1}) + (1-\lambda_1)s_{Lj1}(1-s_{Lj1})][m_{j1}+\delta\sum_k\tilde{s}_{k2}(j)m_{k2}(j)] \\
  &- \alpha N\sum_{h\ne j}[\lambda_1s_{Hh1}s_{Hj1} + (1-\lambda_1)s_{Lh1}s_{Lj1}][m_{h1} + \delta\sum_k\tilde{s}_{k2}(h)m_{k2}(h)] \\
  &+\delta\sum_h\tilde{s}_{h1}\frac{\partial\lambda_2(h)}{\partial p_{j1}}[\pi_{H2}(h) - \pi_{L2}(h)]
\end{align*}

This is much the same as the perfect information case. However, looking at the third line, we can see where the firm incorporates their belief updating. Define the derivative in the third line as such:

\begin{gather*}
  \frac{\partial\lambda_2(h)}{\partial p_{j1}} = 
  \begin{cases}
    0 = \frac{\alpha\lambda_1(1-\lambda_1)}{\tilde{s}_{j1}^2}[s_{Hj1}s_{Lj1}(s_{Lj1}-s_{Hj1})] \ \text{if $h=j$} \\
    0 = \frac{\alpha\lambda_1(1-\lambda_1)}{\tilde{s}_{h1}^2}[s_{Hh1}s_{Lh1}(s_{Lj1}-s_{Hj1})] \ \text{if $h\ne j$} \\
  \end{cases}
\end{gather*}

From this, the firm updates its belief about the distribution of consumer types upon observing period one consumption. The magnitude and direction of this update depends on how likely that choice is for each type. In total, the third line of the derivative interacts the probability of purchasing item $h$ with the effect of price on the firm's posterior belief and the value of its belief in one type or the other.

Now, we can solve for $p_{j1}$ via the following steps. First,

\begin{align*}
  \alpha N\mathbb{A}[p_{j1} - c_j &+ \delta\sum_k\tilde{s}_{k2}(j)m_{k2}(j)] =\\
  &- N\tilde{s}_{j1} + \alpha N\sum_{h\ne j}[\lambda_1s_{Hh1}s_{Hj1} + (1-\lambda_1)s_{Lh1}s_{Lj1}][m_{h1} + \delta\sum_{k}\tilde{s}_{k2}(h)m_{k2}(h)] \\
  &-N\delta\sum_h\tilde{s}_{h1}\frac{\partial\lambda_2(h)}{\partial p_{j1}}[\pi_{H2}(h) - \pi_{L2}(h)]
\end{align*}

Now, we isolate $p_{j1}$:

\begin{align*}
  \alpha N\mathbb{A}p_{j1} &= \alpha N\mathbb{A}c_j - \alpha N\mathbb{A}\delta\sum_k\tilde{s}_{k2}(j)m_{k2}(j) \\
  &- N\tilde{s}_{j1} + \alpha N\sum_{h\ne j}[\lambda_1s_{Hh1}s_{Hj1} + (1-\lambda_1)s_{Lh1}s_{Lj1}][m_{h1}+\delta\sum_k\tilde{s}_{k2}(h)m_{k2}(h)] \\
  &- N\delta\sum_h\tilde{s}_{h1}[\frac{\partial\lambda_2(h)}{\partial p_{j1}}(\pi_{2H}(h) - \pi_{2L}(h))]
\end{align*}

And simplify:

\begin{align*}
  p_{j1} &= c_j -\delta\sum_k\tilde{s}_{k2}(j)m_{k2}(j) - \frac{\tilde{s}_{j1}}{\alpha A} \\
  &+ \frac{\sum_{h\ne j}[\lambda_1s_{Hh1}s_{Hj1} + (1-\lambda_1)s_{Lh1}s_{Lj1}]}{\mathbb{A}}[m_{h1} + \delta\sum_k\tilde{s}_{k2}(h)m_{k2}(h)] \\
  &- \frac{\delta}{\alpha \mathbb{A}}[\sum_h\tilde{s}_{h1}\frac{\partial\lambda_2(h)}{\partial   p_{j1}}(\pi_{2H}(h) - \pi_{2L}(h))] \\
\end{align*}


Deconstructing the optimal price, optimal price is composed of the marginal cost, minus the own-product discounted continuation value, plus a demand markup term and the the value of alternative products, and lastly a term referring to the learning adjustment from period one behavior. In the imperfect-information specification, today's prices affect current demand and the distribution of future states as well as the information that can be learned from consumer purchases. Period one prices alter the likelihood of purchase, thus affecting the posterior over consumer type and, by extension, optimal pricing in period two.

%============================================================%
\section{Supply Simulation}
%============================================================%

In the supply simulation, I set the following parameters

\begin{table}[htbp]
\centering
\caption{Model Parameters}
\label{tab:model_parameters}
\renewcommand{\arraystretch}{1.2}
\begin{tabular}{lll}
\toprule
\textbf{Parameter} & \textbf{Symbol} & \textbf{Assigned Value} \\
\midrule
Base Brand Utilities & $\beta_0$ & $[2.2, 2.0]$ \\
High-Type Switching Utility & $\gamma_H$ & $2.0$ \\
Low-Type Switching Utility & $\gamma_L$ & $0.1$ \\
Price Sensitivity Coefficient & $\alpha$ & $-1.5$ \\
Discount Factor & $\delta$ & $0.95$ \\
Marginal Unit Costs & $c$ & $[1.2, 1.0]$ \\
Number of Inside Goods & $J$ & $2$ \\
Initial State Distribution & $s_0$ & $[0.5, 0.5]$ \\
Price Optimization Bounds & $[P_{\min}, P_{\max}]$ & $[\$1.00, \$8.00]$ \\
\bottomrule
\end{tabular}
\end{table}

Then, I give consumers a prior $j_0$ so period one prices are informed by switching in the imperfect case. In table \ref{tab:value_of_information}, the optimal prices, posterior beliefs, and profits are displayed. Firms lower prices on the lower cost good, directing consumers to it to raise profits. Firms offer higher discounts when the distribution of consumers has less high-switchers, needing to induce low-types to switch via price effects. The firm consistently overestimates the amount of high types in the population by a magnitude of anywhere from 10-15 percent. We see profits maximized in the case akin to perfect information where all individuals are high type. When comparing the case of perfect and imperfect information profits, the value of information is between 1.2-1.9 cents per unit. This cost of information is maximized when the distribution is evenly split, implying the learning adjustment cost from the imperfect cost specification is meaningfully impacting firm profits.

This result points to consumer privacy policies being a viable way to limit tracking a consumer's data. Firms who are unwilling to incur the losses of learning will stop tracking as diligently and the firms large enough to weather the losses will need some other form of enforcement (fines, bans, etc.). In physical goods spaces, this shift to privacy is perhaps less poignant. In digital markets, firms are able to track every action a consumer takes on their domain and can use that information to target advertising, discounts, or otherwise induce prolonged usage of the product.  

%============================================================%
\section{Conclusion}
%============================================================%

In closing, this paper implements a dynamic discrete choice model of consumer demand that incorporates a love of variety and strategic pricing by firms. The model is estimated using scanner data, and counterfactual analyses are conducted to evaluate the effects of different market structures on consumer welfare and market outcomes. The paper contributes to the literature on love of variety, price-discrimination, and dynanmic consumer behavior, by providing empirical evidence on the effects of strategic pricing in a dynamic demand setting. 

I find that consumers experience disutility from switching, needing to be compensated \$2.20 to switch from one flavor to another. I find insignificant pricing effects, even with endogeneity intervention implying a need for a stronger instrumental variable. This simple version of the model suggests that there is significant runway for saturation with other product characteristics, random coefficients, and a richer sample.

The next iteration of this work will first augment the sample to include 2022-2024, as well as add markets. Next, I will nail down the function form of utility to better implement the idea of attribute satiation. There is also room to add more characteristics like calories, grams of protein, grams of sugar, and type of yogurt. A brand effects pathway also makes sense, seeing if households are brand loyal or loyal to flavors. This decomposition of switching behavior would be very useful for the product introduction counterfactual especially. Lastly, the supply side will be an ongoing effort. With that being my primary contribution, it is obviously a priority and thus will be a major focus of forthcoming research.

%============================================================%
\newpage

\bibliography{lov_bib/lov_bib.bib}

\newpage
\section{Tables and Figures}
% ==============================================================================
% TABLE 1: NAIVE SPECIFICATIONS (NO SATIATION / NO GAMMAS)
% ==============================================================================
\begin{table}[htbp]
\centering
\caption{Naive Choice Model Specifications (No Satiation Dynamics)}
\label{tab:naive_specifications}
\renewcommand{\arraystretch}{1.2}
\resizebox{\textwidth}{!}{%
\begin{tabular}{l *{4}{r} c *{4}{r}}
\toprule
& \multicolumn{4}{c}{\textbf{Standard Logit (No Gammas)}} & & \multicolumn{4}{c}{\textbf{Random-Coefficient Logit (No Gammas)}} \\
\cmidrule{2-5} \cmidrule{7-10}
\textbf{Parameter} & \textbf{Estimate} & \textbf{Std. Error} & \textbf{$z$-stat} & \textbf{$p$-value} & & \textbf{Estimate} & \textbf{Std. Error} & \textbf{$z$-stat} & \textbf{$p$-value} \\
\midrule
Constant & 0.9455 & 0.0397 & 23.837 & 0.0000$^{***}$ & & 0.9459 & 0.0397 & 23.850 & 0.0000$^{***}$ \\
$\beta_{\text{berry}}$ (Mean) & -0.7635 & 0.0126 & -60.401 & 0.0000$^{***}$ & & -0.7632 & 0.0127 & -60.298 & 0.0000$^{***}$ \\
$\beta_{\text{plain}}$ (Mean) & -1.3899 & 0.0179 & -77.587 & 0.0000$^{***}$ & & -1.3930 & 0.0218 & -63.784 & 0.0000$^{***}$ \\
SD $\beta_{\text{berry}}$ & \multicolumn{4}{c}{---} & & 0.0008 & 0.0304 & 0.027 & 0.9785 \\
SD $\beta_{\text{plain}}$ & \multicolumn{4}{c}{---} & & 0.0303 & 0.1041 & 0.291 & 0.7711 \\
Price ($\alpha$) & -0.9836 & 0.0299 & -32.878 & 0.0000$^{***}$ & & -0.9836 & 0.0299 & -32.880 & 0.0000$^{***}$ \\
Control Function ($\sigma$) & 0.1670 & 0.0292 & 5.721 & 0.0000$^{***}$ & & 0.1654 & 0.0292 & 5.668 & 0.0000$^{***}$ \\
\midrule
Optimization Success & \multicolumn{4}{c}{True} & & \multicolumn{4}{c}{True} \\
Final Log-Likelihood & \multicolumn{4}{c}{-71,955.73} & & \multicolumn{4}{c}{-71,955.75} \\
\bottomrule
\multicolumn{10}{l}{\footnotesize $^{***}p<0.001$, $^{**}p<0.05$, $^*p<0.10$.}
\end{tabular}%
}
\end{table}

% ==============================================================================
% TABLE 2: WTP FOR NAIVE SPECIFICATIONS
% ==============================================================================
\begin{table}[htbp]
\centering
\caption{Willingness to Pay (WTP in \$) for Naive Specifications}
\label{tab:wtp_naive}
\renewcommand{\arraystretch}{1.2}
\begin{tabular}{l r r}
\toprule
\textbf{Attribute / Parameter} & \textbf{Standard Logit WTP (\$)} & \textbf{Random Coefficient Logit WTP (\$)} \\
\midrule
Berry Preference (vs. Other) & \$0.78 & \$0.78 \\
Plain Preference (vs. Other) & \$1.41 & \$1.42 \\
\bottomrule
\end{tabular}
\end{table}

\vspace{0.5cm}

% ==============================================================================
% TABLE 3: FLAVOR SATIATION SPECIFICATIONS (WITH GAMMAS)
% ==============================================================================
\begin{table}[htbp]
\centering
\caption{Flavor-Specific Satiation Choice Model Specifications}
\label{tab:satiation_specifications}
\renewcommand{\arraystretch}{1.2}
\resizebox{\textwidth}{!}{%
\begin{tabular}{l *{4}{r} c *{4}{r}}
\toprule
& \multicolumn{4}{c}{\textbf{Standard Flavor Satiation Logit}} & & \multicolumn{4}{c}{\textbf{Random-Coefficient Flavor Satiation Logit}} \\
\cmidrule{2-5} \cmidrule{7-10}
\textbf{Parameter} & \textbf{Estimate} & \textbf{Std. Error} & \textbf{$z$-stat} & \textbf{$p$-value} & & \textbf{Estimate} & \textbf{Std. Error} & \textbf{$z$-stat} & \textbf{$p$-value} \\
\midrule
Constant & 0.9447 & 0.0397 & 23.824 & 0.0000$^{***}$ & & 0.9353 & 0.0398 & 23.477 & 0.0000$^{***}$ \\
$\beta_{\text{berry}}$ (Mean) & -0.7611 & 0.0126 & -60.171 & 0.0000$^{***}$ & & -0.7500 & 0.0126 & -59.687 & 0.0000$^{***}$ \\
$\beta_{\text{plain}}$ (Mean) & -1.3894 & 0.0179 & -77.464 & 0.0000$^{***}$ & & -1.3994 & 0.0155 & -90.252 & 0.0000$^{***}$ \\
$\gamma_{\text{berry}}$ (Mean) & -0.0678 & 0.0221 & -3.065 & 0.0022$^{**}$ & & -10.0000 & 0.3606 & -27.730 & 0.0000$^{***}$ \\
$\gamma_{\text{plain}}$ (Mean) & -0.0130 & 0.0177 & -0.735 & 0.4625 & & -1.7404 & 0.5154 & -3.377 & 0.0007$^{***}$ \\
SD $\beta_{\text{berry}}$ & \multicolumn{4}{c}{---} & & 0.0000 & 0.0000 & $\infty$ & 0.0000$^{***}$ \\
SD $\beta_{\text{plain}}$ & \multicolumn{4}{c}{---} & & 0.0000 & 0.0000 & $\infty$ & 0.0000$^{***}$ \\
SD $\gamma_{\text{berry}}$ & \multicolumn{4}{c}{---} & & 6.9874 & 0.2905 & 24.053 & 0.0000$^{***}$ \\
SD $\gamma_{\text{plain}}$ & \multicolumn{4}{c}{---} & & 4.4204 & 1.2400 & 3.565 & 0.0004$^{***}$ \\
Price ($\alpha$) & -0.9830 & 0.0299 & -32.868 & 0.0000$^{***}$ & & -0.9759 & 0.0301 & -32.468 & 0.0000$^{***}$ \\
Control Function ($\sigma$) & 0.1671 & 0.0292 & 5.726 & 0.0000$^{***}$ & & 0.1762 & 0.0291 & 6.045 & 0.0000$^{***}$ \\
\midrule
Optimization Success & \multicolumn{4}{c}{True} & & \multicolumn{4}{c}{True} \\
Final Log-Likelihood & \multicolumn{4}{c}{-71,947.72} & & \multicolumn{4}{c}{-71,894.17} \\
\bottomrule
\multicolumn{10}{l}{\footnotesize $^{***}p<0.001$, $^{**}p<0.05$, $^*p<0.10$.}
\end{tabular}%
}
\end{table}

% ==============================================================================
% TABLE 4: WTP FOR SATIATION SPECIFICATIONS
% ==============================================================================
\begin{table}[htbp]
\centering
\caption{Willingness to Pay (WTP in \$) for Satiation Specifications}
\label{tab:wtp_satiation}
\renewcommand{\arraystretch}{1.2}
\begin{tabular}{l r r}
\toprule
\textbf{Attribute / Parameter} & \textbf{Standard Flavor Satiation WTP (\$)} & \textbf{RC Flavor Satiation WTP (\$)} \\
\midrule
Berry Preference (vs. Other) & \$0.77 & \$0.77 \\
Plain Preference (vs. Other) & \$1.41 & \$1.43 \\
Satiation Disutility (Berry) & \$1.41 & \$1.43 \\
Satiation Disutility (Plain) & \$0.07 & \$10.25 \\
\bottomrule
\end{tabular}
\end{table}

% ==============================================================================
% TABLE 5: OPTIMAL PRICING
% ==============================================================================
\begin{table}[htbp]
\centering
\caption{Information Structure, Optimal Pricing, and Value of Information Across Population Mixes}
\label{tab:value_of_information}
\renewcommand{\arraystretch}{1.25} % Adds comfortable vertical breathing room
\resizebox{\textwidth}{!}{% Fits perfectly to text margins without crowding
\begin{tabular}{l c c r r r r r r r}
\toprule
\textbf{Population Mix} & \textbf{\% High} & \textbf{\% Low} & \multicolumn{2}{c}{\textbf{Imperfect Price ($P_1$)}} & \multicolumn{2}{c}{\textbf{Posterior Beliefs}} & \textbf{Imperfect} & \textbf{Perfect} & \textbf{Value of} \\
& $\left(\frac{N_H}{N}\right)$ & $\left(\frac{N_L}{N}\right)$ & \textbf{Good 1} & \textbf{Good 2} & \textbf{G1} & \textbf{G2} & \textbf{Profit} & \textbf{Profit} & \textbf{Info} \\
\midrule
100\% LOW          & 0\%   & 100\% & 2.0633 & 1.8646 & 0.0000 & 0.0000 & 0.6031 & 0.6031 & -0.0000 \\
25\% HIGH, 75\% LOW & 25\%  & 75\%  & 2.1847 & 1.9944 & 0.3727 & 0.3641 & 0.7820 & 0.7999 &  0.0178 \\
50\% HIGH, 50\% LOW & 50\%  & 50\%  & 2.2819 & 2.1049 & 0.6510 & 0.6448 & 0.9773 & 0.9966 &  0.0194 \\
75\% HIGH, 25\% LOW & 75\%  & 25\%  & 2.3542 & 2.1887 & 0.8528 & 0.8504 & 1.1811 & 1.1934 &  0.0123 \\
100\% HIGH         & 100\% & 0\%   & 2.4067 & 2.2494 & 1.0000 & 1.0000 & 1.3902 & 1.3902 &  0.0000 \\
\bottomrule
\end{tabular}%
}
\end{table}

\end{document}
